\documentclass[10pt,a4paper]{amsart}
\usepackage[margin=19mm]{geometry}
\usepackage{amsmath,amssymb,mathtools}
\usepackage[scr=rsfs,cal=euler]{mathalfa}
\usepackage{xfrac}
\usepackage[unicode,hidelinks,bookmarks=false]{hyperref}
\newcommand{\ie}{i.e.\ }
\newcommand{\cf}{cf.\ }
\newcommand{\resp}{respectively\ }
\newcommand{\Subsection}{\subsection}
\providecommand{\nobreakdash}{\nobreak\hbox{-}}
\setlength{\emergencystretch}{2em}
\multlinegap0pt
\usepackage{amssymb}
\usepackage{dsfont}
\usepackage{float}
\usepackage{mathtools}
\usepackage{xcolor}
\newtheorem{lemma}{Lemma}
\newcommand{\trm}{\mathtt t^-}
\newcommand{\trp}{\mathtt t^+}
\title{Asymptotic behavior of Carleman weight functions and application to controllability}
\author{Ariel A. P\'erez}
\date{Source: arXiv:2412.19892v2 (CC BY 4.0)}
\begin{document}
\maketitle

\noindent\emph{Source and licence.} Asymptotic behavior of Carleman weight functions and application to controllability. Version arXiv:2412.19892v2 (CC BY 4.0), distributed by arXiv under the Creative Commons Attribution 4.0 International licence, http://creativecommons.org/licenses/by/4.0/, as recorded in the arXiv licence metadata for this exact version and shown on the abstract page https://arxiv.org/abs/2412.19892v2. Full licence notice: \texttt{licenses/arxiv-2412.19892-CC-BY-4.0.txt}.\par\smallskip

\noindent\textit{Editorial scope.} Counted unit: the article's lemma lem:time:discrete:weight (source0.tex lines 548--555). The article's complete original proof of it is delivered with it (source0.tex lines 556--591, one \texttt{proof} environment of 36 lines). No mathematics is written, repaired or re-derived: the statement and the proof are the article's own lines, in the article's own notation, and no retained line is substituted or re-flowed.  No further source-local context is needed: every cross-reference of every retained line resolves inside the two delivered spans.  Nothing was omitted from any retained span, so the package carries no omission record.  No retained line contains a citation, so the wrapper carries no bibliography and no bibliography entry.  No retained line contains a figure environment, an image-inclusion command or drawing code, so no image is delivered and the package declares no asset dependency.  Editorial formatting only, in addition: an AMS \texttt{amsart} preamble that restates the article's own declarations and macros used by the retained lines, namely the environments \texttt{lemma} on separate counters, together with the standard packages those lines need.  The retained environments print in this document as Lemma 1 in order of appearance, whereas the article prints the same items with the numbering of its own full document.  The complete native source of the article as distributed in the arXiv e-print for this exact version is bundled unchanged as \texttt{sources/170806/source0.tex}.
\begin{lemma}\label{lem:time:discrete:weight}
    Provided $\Delta t\tau(T^{3}\delta^{2})^{-1}\leq 1$, we have
    \begin{equation}\label{eq:theta:D:t}
        \trm(\theta^{p}e^{s\varphi})D_{t}(\theta^{-p}e^{-s\varphi})=\trm(\theta^{p}e^{s\varphi}\partial_{t}(\theta^{-p}e^{-s\varphi}))+\mathcal{E}(\Delta t) \mathcal{O}_{\lambda}(1),
    \end{equation}
    where $\mathcal{E}(\Delta t):=\mathcal{E}(\Delta t, p,\delta,\tau,T):=p(p+1)\frac{\Delta t}{\delta^{p+4}T^{2}}+p\frac{\Delta t}{\delta^{p+3}T^{2}}
    +2p\tau\frac{\Delta t}{\delta^{p+4}T^{4}}+\tau^{2} \frac{\Delta t}{\delta^{p+4}T^{6}}+\tau\frac{\Delta t}{\delta^{p+3}T^{4}}$.
\end{lemma}

\begin{proof}
We use the first-order Taylor formula for  $f:=\theta^{-p}e^{-s\varphi}$ to write
\begin{equation}\label{eq:time:taylor}
    \trp D_{t}(f)=\partial_{t}f+\Delta t\int_{0}^{1}(1-\gamma)\partial_{t}^{2}f(t+\gamma \Delta t)d\gamma.
\end{equation}
Let us focus on the integral remainder of the above expression. We have
\begin{equation*}
\begin{aligned}
    \partial_{t}^{2}f=&(p(p+1)\theta^{-p-2}(\partial_{t}\theta)^{2}-p\theta^{-p-1}\partial_{t}^{2}\theta)e^{-s\varphi}+2\tau\varphi d\theta^{-p-1}(\partial_{t}\theta)^{2}e^{-s\varphi}\\
    &+\theta^{-p}\left[ \tau^{2}\varphi^{2}(\partial_{t}\theta)^{2}e^{-s\varphi}-\tau\partial_{t}^{2}\theta\varphi e^{-s\varphi}\right].
      \end{aligned}
\end{equation*}
We notice that there exists a constant $C>0$, only depending on $p$, such that
\begin{equation}\label{ine:theta:d:-d}
    \max_{t\in [0,T]} \theta^{p}(t)\theta^{-p}(t+\gamma\Delta t)\leq \frac{C}{\delta^{p}}.
\end{equation} Indeed, for $p\in\{0,1,2,\ldots\}$ it follows that
\begin{equation}\label{ine:max:theta}
    \max_{t\in[0,T]}\theta^{-p}(t+\gamma\Delta t)\leq   \frac{T^{2p}(1+2\delta)^{2p}}{4^{p}}\leq T^{2p},
    \end{equation}
    and 
    \begin{equation}\label{ine:theta:d}
    \max_{t\in[0,T]} \theta^{p}(t)\leq \frac{C}{\delta^{p}T^{2p}}.
\end{equation} 
So \eqref{ine:theta:d:-d} holds. And we also have
\begin{equation}\label{ine:theta:d:2}
    \max_{t\in[0,T]}\theta^{(j)}(t)\leq \frac{1}{\delta^{j+1}T^{j+2}},\quad j=0,1,\ldots
\end{equation}
Thus, multiplying \eqref{eq:time:taylor} by $\theta^{p}e^{s\varphi}$, and then applying \eqref{ine:theta:d:-d}-\eqref{ine:theta:d:2} to the result we obtain
\begin{equation}
\begin{aligned}
    \theta^{p}e^{s\varphi}\trp D_{t}(\theta^{-p}e^{-s\varphi})=&\theta^{p}e^{s\varphi}\partial_{t}(\theta^{-p}e^{-s\varphi})+p(p+1)\frac{\Delta t}{\delta^{p+4}T^{2}}\mathcal{O}_{\lambda}(1)+p\frac{\Delta t}{\delta^{p+3}T^{2}}\mathcal{O}_{\lambda}(1)\\
    &+2p\tau\frac{\Delta t}{\delta^{p+4}T^{4}}\mathcal{O}_{\lambda}(1)+\tau^{2} \frac{\Delta t}{\delta^{p+4}T^{6}}\mathcal{O}_{\lambda}(1)+\tau\frac{\Delta t}{\delta^{p+3}T^{4}}\mathcal{O}_{\lambda}(1),
\end{aligned}
\end{equation}
where the conditions $\Delta t\tau(\delta^{2} T^{3})^{-1}\leq 1$ and $\|\varphi\|_{\infty}=\mathcal{O}_{\lambda}(1)$ are used to control the exponential term $e^{\tau(\theta(t)-\theta(t+\gamma\Delta t)\varphi)}=\mathcal{O}_{\lambda}(1)$. Then, by shifting the above expression using $\trm$ the estimate \eqref{eq:theta:D:t} holds.
\end{proof}

\end{document}
